\section{Parameter choices and completion of the proof}\label{sec:parameters}

Recall that \(x=N^{-1/6}\) and
\(x^{12}\le\lambda\le x^{p_*}\).  The parameters below are fixed
independently of \(N\).  We first choose the savings used to control
size changes, then the final-strength exponent \(p_*\), and only then
the truncation and loss parameters for direct changes. The essential
order is that the size Taylor degrees determine the polynomial loss
constants before the cutoffs and \(p_*\) are selected. Direct Taylor
degrees may be larger, because their estimates do not enter that size
choice. Every degree, moment order, and loss exponent below is fixed
before \(N\to\infty\).

\begin{enumerate}
\item Fix \(T>1\) and the neighborhoods, tapers, and constants
\(c_1,c_b\) in \eqref{eq:P1}--\eqref{eq:P7}, leaving strict room in
the admissibility and derivative bounds.  The global saving \(c_G\)
in \eqref{eq:G1} is available before choosing \(s_*\) near \(1/2\)
or the ultimate small \(p_*\).  In its proof, the coarse clock estimate
\eqref{eq:G3} supplies a fixed saving first.  Choose \(a_g\) below
that saving, then \(d_g\) small enough to absorb the fixed
inverse-kernel loss in the covariance estimate following
\eqref{eq:G4}, and finally the PE power slack
for this global comparison small relative to \(d_g\).  The endpoint
condition \eqref{eq:P9} uses the common base temperature as explained
in Section~\ref{sec:global}.  It requires \(\lambda\to0\), without
requiring \(\lambda\le x^{a_g}\).  Thus later decreases of \(p_*\)
do not decrease \(c_G\).

\item Set \(k_c=0.03\).  Choose
\(\tau\ll\min(10^{-4},c_G)\), put \(s_*=1/2-\tau\), and then
choose \(\eta\ll\tau\) and \(\tau_0\ll\eta\), with
\(s_0=s_*-\tau_0\), as in \eqref{eq:C1}--\eqref{eq:C2}.
Choose \(b_s>0\) so that
\[
 2(\tau+\tau_0)+b_s<\min(c_G/12,k_c/3),\qquad
 b_s<2(\eta-\tau_0),
\]
and then \(b_\ell\ll b_s\).  These inequalities leave room for the
local comparison errors throughout \(\lambda\ge x^{12}\).
Choose also \(h_e\ll\min(1,c_G)\).

Both final-strength estimates in \eqref{eq:C2} now have a fixed
saving \(c>0\) for every sufficiently small \(p_*\).  Indeed, with
\(\lambda=x^{p_*}\), \(U=x\lambda^{-s_*/2}\), and the edge seed
off, the differentiated bound in its proof has powers
\[
 4+c_G-p_*,\qquad 5+p_*(2\eta-5s_*/2),\qquad 6-2p_*.
\]
All remain separated above \(4\) as \(p_*\) decreases.  The
smoothing error \(O(x^{5+c'})\) in \eqref{eq:P11}, divided by
\(\lambda\), has exponent \(5+c'-p_*>4\) as well.
Likewise \eqref{eq:D11} supplies a fixed
\(\delta_*>0\), for example \((1/200)/50\), under the preliminary
restriction \(p_*<(1/200)/100\).  Neither saving depends on the
remaining choice of \(p_*\).

\item Fix the size row and pressure Taylor degrees using
\(\delta_*\).  In the covariance equation proving \eqref{eq:D15},
choose the row degree so that the Taylor remainder before inversion,
apart from the \(O(m^{-1})\) contacts, is smaller than \(x^6\) by
a fixed power.  This is possible by \eqref{eq:D11} and requires no
further information about \(p_*\).  For these fixed degrees the
inverse has fixed polynomial losses in
\((\lambda V_i^2)^{-1}\) and \(1+M\); the low sources have a fixed
loss \(\lambda^{-C}\).  Fix constants bounding these losses before
choosing the cutoffs, large enough to cover all the estimates below.

Now choose \(a_0\) small relative to \(c,\delta_*\), take \(b_H\)
to be a sufficiently small fixed fraction of \(a_0\), then choose
\(a_1\) small relative to \(a_0,b_H\), and \(a_2\) small relative
to \(a_1\).  Finally choose \(p_*\) small relative to these gaps and
to \(h_e\), including the restriction needed for the cap on \(Y\)
after \eqref{eq:D11}.  With \(V_i=x^{a_i}\), the high-mass terms
of \eqref{eq:E3} then have the following positive margins over
\(x^5\):
\[
\begin{gathered}
 1-2C_Ha_1>0,\qquad b_H-C_Ha_1>0,\qquad b_H-Cp_*>0,\\
 a_1-C_Ha_2-p_*s_*>0,\qquad
 \delta_*/2-C_Ha_2-p_*s_*>0.
\end{gathered}
\]
The first two margins cover two high stochastic factors and a high
stochastic factor paired with a high deterministic factor.  The third
covers a high deterministic factor times the low-source estimate
\eqref{eq:D16}.  The last two cover a sole high stochastic factor,
using respectively the parity-weighted integral \(V_1Y^2\) or a
further factor \(x^{\delta_*/2}Y^2\).  In both cases
\(Y^2=x^2\lambda^{-s_*}Z^2\).  The size projection errors
\(x^6\lambda^{-C}=x^{6-Cp_*}\) also have the positive margin
\(1-Cp_*\).  No subsequent step decreases \(p_*\).

\item Fix the number of strengthening stages in \eqref{eq:C11}.
Choose \(\zeta\) small relative to the power margins in
Section~\ref{sec:budgets}; in particular, it is smaller than \(p_*\)
times each required gap in a \(\lambda\)-exponent, including the gaps
involving \(b_\ell\) and \(\tau_0\).  Then choose the local PE
power slacks small relative to \(\zeta\), and the buffer exponent
\(\rho\) small enough for \eqref{eq:C18}.  The remaining finite
profile and moment iterations, and their subpower losses, can be
chosen with these parameters fixed.

The extra optimized excess
\(O(x^{4-2\zeta}u\Lambda^{1/3})\) fits the budget at
\eqref{eq:C3}: after division by \(\lambda u\), its ratio to
\(x^4\lambda^{-2s_0}\) is
\(x^{-2\zeta}\lambda^{k_c/3-1+2s_0}\), which is power small by
these choices.  The final-strength switch error
\(O(x^{5+c'})\) from \eqref{eq:P11} also fits after division by
\(\lambda u\).  At the smallest dyad \(u\asymp U\), its ratio
to \(U^4\) is at most
\(x^{c'}\lambda^{5s_*/2-1}\); larger dyads improve this bound.
Both exponents are positive, so the switch leaves room for the
remaining losses.

\item For the direct paths, choose the coarse cutoff exponent
\(\kappa>0\) small relative to the gains in
Lemma~\ref{lem:statistics-coarse-switch}.  More precisely, its
large-\(Y\) event costs \(x^{-4\kappa}\), so choose
\[
 4\kappa<\min\left\{p_*(1-2s_*),\,
                       (1-\eta)/s_*-2\right\}.
\]
Both quantities on the right are positive by the choices of
\(s_*,\eta\).  The gains used in the
remaining direct expansion are
\[
 1-2s_*=2\tau>0,\qquad
 (1-\eta)/2-s_*=\tau-\eta/2>0,\qquad
 s_*/2-\eta>0.
\]
A factor \(\lambda^g\) with any of these exponents \(g>0\) is at
most \(x^{p_*g}\).  Choose \(\delta'>0\) smaller than \(\kappa\)
and small enough that
\[
 C_0\delta'<p_*\bigl((1-\eta)/2-s_*\bigr),
\]
with strict room for the later losses.  Thus the direct Taylor error
\(x^{5-C_0\delta'}\lambda^{(1-\eta)/2-s_*}\) saves a power over
\(x^5\).  Now fix the direct pressure Taylor degree large enough
also for this \(\delta'\).  The initial direct projection cost
satisfies
\(x^6\lambda^{-\eta}\le x^{6-12\eta}\), so its margin over
\(x^5\) is \(1-12\eta>0\).

Choose the fixed moment and change-of-measure orders large enough
that truncation and transfer lose only small fractions of the margins
just retained.  Then choose the geometric loss in \eqref{eq:D1} and
\eqref{eq:D10}, and the subpower losses, still smaller.  The smooth
density comparisons on intervals of length \(O(b_\epsilon)\) may
require higher fixed orders after this geometric choice; they change
constants only.  Finally take the stencil order large enough for
both the fixed size pressure polynomial and the site-weight smoothing
in Section~\ref{sec:reduction}.  Enlarging this order does not require
enlarging the size Taylor degree.
\end{enumerate}

These choices give a common positive saving in the stopped increment
bounds.  For a direct strength sweep, Lemma~\ref{lem:statistics-coarse-switch}
handles \(Y>x^\kappa\); on the complement apply the allocation
and compression estimates \eqref{eq:E2}, \eqref{eq:E4} to the
pressure polynomial \eqref{eq:E3}. The two multiplier sweeps use the
separate switch bounds of the same lemma. The size sweep uses the
capped low statistics, \eqref{eq:D15}, and the size versions of these
polynomial estimates.  Integrating the direct strength
parameter costs only \(O(\log(1/x))\).  The size segment has bounded
length, and the multiplier bound has at most the integrable endpoint
factor \(\chi_e^{-1/2}\).  The fixed-degree dyadic and projection-cut
sums cost subpower factors.  Decreasing the common saving absorbs
these integration and summation costs.

We spell out the uniformity needed when the site count varies.  At
each fixed size-path point \(\theta\), every estimate above is
applied to the particular integer \(m\) being tested, uniformly for
\(N/2\le m\le3N\).  In particular, the averaging in
\eqref{eq:C1}--\eqref{eq:C2} is over \(\gamma\) at this fixed
\((\theta,m)\).  The physical clocks are common to these site
counts.  Their absolute dotted-clock norms are bounded up to subpower
loss as explained after \eqref{eq:P8}; integrating any remaining
dotted weight in mass therefore gives a common integrable envelope
in \(\theta\).  Thus the bound is uniform in \(m\) before
integrating \(\theta\).

Let \(w_m(\theta)\ge0\) be measurable probability weights on the
allowed integers.  Multiplying the parameter-averaged bound at
\((\theta,m)\) by \(w_m(\theta)\) and summing preserves its
common envelope, since \(\sum_m w_m(\theta)=1\).  A measurable
choice of a single \(m(\theta)\) is the special case of a point
mass.  Fixed-size path approximation makes the pressure tests and
statistics measurable, and absolute continuity supplies their a.e.
path derivatives.  Fubini first removes a null set of \(\theta\)
for each parameter-averaged estimate.  At fixed \(N\) there are only
finitely many allowed \(m\), so their exceptional sets can then be
united into one null set before weighting and integrating.

For each summand the change of measure is used only on the central
stop \(\mathcal K_m(s)\le3x^\varepsilon\).  The comparisons with
finitely many added rows in \eqref{eq:A6}--\eqref{eq:A9} have bounded
covariance changes and estimate their remainders with moments of the
original \(m\)-site system.  They therefore require this same stop,
without a new stopping condition for the added-row systems.  Fix the
common saving before taking the small Laplace exponent
\(\varepsilon\) in Section~\ref{sec:reduction}; this exponent is
distinct from the geometric slack. The density bounds on the stop hold
for every fixed finite order, so they cover all transfer orders chosen
above, for every sufficiently small fixed \(\varepsilon>0\) and each
fixed test coefficient \(q>0\) in \(s=qx^{1+\varepsilon}\).

For one such test, the same parameter bounds cover the finitely many
segments in the comparison from \(N\) to any \(N'\in[N,2N]\).
Summing their stopped bounds and then choosing \(\gamma\) gives the
single route used in Section~\ref{sec:reduction}. Different tests may
choose different routes: their bare endpoint laws are independent of
\(\gamma\). Thus the later interpolation at finitely many Laplace
tests requires no parameter choice uniform over an infinite test family.

This proves Proposition~\ref{prop:increment}, including its version
with a measurably varying site count or site weights.
Section~\ref{sec:reduction} then proves Theorem~\ref{thm:moments}
and proves Proposition~\ref{def:law}.  Removing the independent
Gaussian completion and dividing by the exact finite-size standard
deviation, as done there, proves Theorem~\ref{thm:main} for every
fixed \(\beta>1\).
